\section{USAMO 2015}
        \begin{enumerate}
        	\item (\textit{USAMO 2015}) Solve in integers the equation

 \[x^2+xy+y^2 = \left(\frac{x+y}{3}+1\right)^3.\]


 



\textbf{Solution}

We first notice that both sides must be integers, so $\frac{x+y}{3}$ must be an integer. 

 We can therefore perform the substitution $x+y = 3t$ where $t$ is an integer. 

 Then:

 $(3t)^2 - xy = (t+1)^3$

 $9t^2 + x (x - 3t) = t^3 + 3t^2 + 3t + 1$

 $4x^2 - 12xt + 9t^2 = 4t^3 - 15t^2 + 12t + 4$

 $(2x - 3t)^2 = (t - 2)^2(4t + 1)$

 $4t+1$ is therefore the square of an odd integer and can be replaced with $(2n+1)^2 = 4n^2 + 4n +1$

 By substituting using $t = n^2 + n$ we get:

 $(2x - 3n^2 - 3n)^2 = [(n^2 + n - 2)(2n+1)]^2$

 $2x - 3n^2 - 3n = \pm (2n^3 + 3n^2 -3n -2)$

 $x = n^3 + 3n^2 - 1$ or $x = -n^3 + 3n + 1$

 Using substitution we get the solutions: $(n^3 + 3n^2 - 1, -n^3 + 3n + 1) \cup (-n^3 + 3n + 1, n^3 + 3n^2 - 1)$

 



\textbf{Solution 2}

Let $n = \frac{x+y}{3}$.Thus, $x+y = 3n$.We have
 \[x^2+xy+y^2 = \left(\frac{x+y}{3}+1\right)^3 \implies (x+y)^2 - xy = \left(\frac{x+y}{3}+1\right)^3\]Substituting $n$ for $\frac{x+y}{3}$, we have
 \[9n^2 - x(3n-x) = (n+1)^3\]Treating $x$ as a variable and $n$ as a constant, we have
 \[9n^2 - 3nx + x^2 = (n+1)^3,\]which turns into
 \[x^2 - 3nx + (9n^2 - (n+1)^3) = 0,\]a quadratic equation.By the quadratic formula, 
 \[x = \frac{1}{2} \left(3n \pm \sqrt{9n^2 - 4(9n^2 - (n+1)^3)} \right)\]which simplifies to
 \[x = \frac{1}{2} \left(3n \pm \sqrt{4(n+1)^3 - 27n^2} \right)\]Since we want $x$ and $y$ to be integers, we need $4(n+1)^3 - 27n^2$ to be a perfect square.We can factor the aforementioned equation to be 
 \[(n-2)^2 (4n+1) = k^2\]for an integer $k$. Since $(n-2)^2$ is always a perfect square, for $(n-2)^2 (4n+1)$ to be a perfect square, $4n + 1$ has to be a perfect square as well.Since $4n + 1$ is odd, the square root of the aforementioned equation must be odd as well.Thus, we have $4n + 1 = a^2$ for some odd $a$.Thus, 
 \[n = \frac{a^2 - 1}{4},\] in which by difference of squares it is easy to see that all the possible values for $n$ are just $n = p(p-1)$, where $p$ is a positive integer. Thus, 
 \[x+y = 3n = 3p(p-1).\]Thus, the general form for 
 \[x = \frac{1}{2} \left(3p(p-1) \pm \sqrt{4(p(p-1)+1)^3 - 27(p(p-1))^2} \right)\] for a positive integer $p$.(This is an integer since $4(p(p-1)+1)^3 - 27(p(p-1))^2$ is an even perfect square (since $4(p(p-1)+1)$ is always even, as well as $27(p(p-1))^2$ being always even) as established, and $3p(p-1)$ is always even as well. Thus, the whole numerator is even, which makes the quantity of that divided by $2$ always an integer.)Since $y = 3n - x$, the general form for $y$ is just
 \[y = 3p(p-1) - \frac{1}{2} \left(3p(p-1) \pm \sqrt{4(p(p-1)+1)^3 - 27(p(p-1))^2} \right)\](This is an integer since $4(p(p-1)+1)^3 - 27(p(p-1))^2$ is an even perfect square (since $4(p(p-1)+1)$ is always even, as well as $27(p(p-1))^2$ being always even) as established, and $3p(p-1)$ is always even as well. Thus, the whole numerator is even, which makes the quantity of that divided by $2$ always an integer, which thus trivially makes 
 \[3p(p-1) - \frac{1}{2} \left(3p(p-1) \pm \sqrt{4(p(p-1)+1)^3 - 27(p(p-1))^2} \right)\] an integer.)for a positive integer $p$.Thus, our general in integers $(x, y)$ is 
 \[(\frac{1}{2} \left(3p(p-1) \pm \sqrt{4(p(p-1)+1)^3 - 27(p(p-1))^2} \right), 3p(p-1) - \frac{1}{2} \left(3p(p-1) \pm \sqrt{4(p(p-1)+1)^3 - 27(p(p-1))^2} \right).\]$\boxed{}$

 -fidgetboss\_4000

 


	\item (\textit{USAMO 2015}) Quadrilateral $APBQ$ is inscribed in circle $\omega$ with $\angle P = \angle Q = 90^{\circ}$ and $AP = AQ < BP$. Let $X$ be a variable point on segment $\overline{PQ}$. Line $AX$ meets $\omega$ again at $S$ (other than $A$). Point $T$ lies on arc $AQB$ of $\omega$ such that $\overline{XT}$ is perpendicular to $\overline{AX}$. Let $M$ denote the midpoint of chord $\overline{ST}$. As $X$ varies on segment $\overline{PQ}$, show that $M$ moves along a circle.


 



\textbf{Solution 1}

We claim the $M$ lies on the circle centered at the midpoint of $\overline{AO}$ that passes through $P$. 

 Proceed via complex numbers. $a=1, b=-1$ and $s, t$ are free parameters. Then, 
 \[x = \frac{1}{2} (s+1+t-s\overline{t}) \text{ and } \Re (p) = \Re (x) \implies p+\overline{p}=x+\overline{x}\] Putting it together, 
 \[\left| \frac{1}{2}-\left(\frac{s+t}{2}\right) \right| = \left| \frac{1}{2} - p \right| \iff (s+t-1)(\overline{s}+\overline{t}-1) = 5-2(p+\overline{p}) = 5 - 2(x+\overline{x})\] 
 \[\iff 3+s\overline{t}+\overline{s}t-s-t-\overline{s}-\overline{t}=5-(s+1+t-s\overline{t}+\overline{s}+1+\overline{t}-\overline{s}t)\] which is clear by expanding. Thus, the proof is complete.

 -sillybone

 



\textbf{Solution 2}

We will use coordinate geometry.

 Without loss of generality,let the circle be the unit circle centered at the origin, 
 \[A=(1,0) P=(1-a,b), Q=(1-a,-b)\],where $(1-a)^2+b^2=1$.

 Let angle $\angle XAB=A$, which is an acute angle, $\tan{A}=t$, then $X=(1-a,at)$.

 Angle $\angle BOS=2A$, $S=(-\cos(2A),\sin(2A))$.Let $M=(u,v)$, then $T=(2u+\cos(2A), 2v-\sin(2A))$.

 The condition $TX \perp AX$ yields: $(2v-\sin(2A)-at)/(2u+\cos(2A)+a-1)=\cot A.$     (E1)

 Use identities $(\cos A)^2=1/(1+t^2)$,  $\cos(2A)=2(\cos A)^2-1= 2/(1+t^2) -1$, $\sin(2A)=2\sin A\cos A=2t^2/(1+t^2)$, we obtain $2vt-at^2=2u+a$.   (E1')

 The condition that $T$ is on the circle yields $(2u+\cos(2A))^2+ (2v-\sin(2A))^2=1$, namely $v\sin(2A)-u\cos(2A)=u^2+v^2$.   (E2)

 $M$ is the mid-point on the hypotenuse of triangle $STX$, hence $MS=MX$, yielding $(u+\cos(2A))^2+(v-\sin(2A))^2=(u+a-1)^2+(v-at)^2$.   (E3)

 Expand (E3), using (E2) to replace $2(v\sin(2A)-u\cos(2A))$ with $2(u^2+v^2)$, and using (E1') to replace $a(-2vt+at^2)$ with $-a(2u+a)$, and we obtain$u^2-u-a+v^2=0$, namely $(u-\frac{1}{2})^2+v^2=a+\frac{1}{4}$, which is a circle centered at $(\frac{1}{2},0)$ with radius $r=\sqrt{a+\frac{1}{4}}$.

 


 



\textbf{Solution 3}

Let the midpoint of $AO$ be $K$. We claim that $M$ moves along a circle with radius $KP$.

 We will show that $KM^2 = KP^2$, which implies that $KM = KP$, and as $KP$ is fixed, this implies the claim.

 $KM^2 = \frac{AM^2+OM^2}{2}-\frac{AO^2}{4}$ by the median formula on $\triangle AMO$.

 $KP^2 = \frac{AP^2+OP^2}{2}-\frac{AO^2}{4}$ by the median formula on $\triangle APO$.

 $KM^2-KP^2 = \frac{1}{2}(AM^2+OM^2-AP^2-OP^2)$.

 As $OP = OT$, $OP^2-OM^2 = MT^2$ from right triangle $OMT$. $(1)$

 By $(1)$, $KM^2-KP^2 = \frac{1}{2}(AM^2-MT^2-AP^2)$.

 Since $M$ is the circumcenter of $\triangle XTS$, and $MT$ is the circumradius, the expression $AM^2-MT^2$ is the power of point $A$ with respect to $(XTS)$. However, as $AX*AS$ is also the power of point $A$ with respect to $(XTS)$, this implies that $AM^2-MT^2=AX*AS$. $(2)$

 By $(2)$, $KM^2-KP^2 = \frac{1}{2}(AX*AS-AP^2)$

 Finally, $\triangle APX \sim \triangle ASP$ by AA similarity ($\angle XAP = \angle SAP$ and $\angle APX = \angle AQP = \angle ASP$), so $AX*AS = AP^2$. $(3)$

 By $(3)$, $KM^2-KP^2=0$, so $KM^2=KP^2$, as desired. $QED$

 



\textbf{Solution 4(synthetic)}

To begin with, we connect $\overline{AT}$ and we construct the nine-point circle of $\triangle AST$ centered at $N_9$.

 Lemma $1$: $AX \cdot AS = AP^2$.We proceed on a directed angle chase. We get $\measuredangle ASP = \measuredangle AQP = \measuredangle QPA$, so $\triangle PAS \sim \triangle XAP$ and the desired result follows by side length ratios.

 Lemma $2$: The locus of $N_9$ as $X$ moves along $\overline{PQ}$ is a circle centered about $A$.We add the midpoint of $\overline{AS}$, $N$, and let the circumradius of $\triangle AST$ be $R$. Taking the power of $A$ with respect to $(N_9)$, we get 
 \[AN_9^2 - \left(\frac{1}{2} R\right)^2 = \text{Pow}_{(N_9)} A = AX \cdot AN = \frac{1}{2} AX \cdot AS = \frac{1}{2} AP^2.\]Hence, $AN_9 = \sqrt{\frac{1}{4}R^2 + \frac{1}{2}AP^2}$, which remains constant as $X$ moves.

 Next, consider the homothety of scale factor $\frac{2}{3}$ about $O$ mapping $N_9$ to $G$. This means that the locus of $G$ is a circle as well.

 Finally, we take a homothety of scale factor $\frac{3}{2}$ about $A$ mapping $G$ to $M$. Hence, the locus of $M$  is a circle, as desired. - Spacesam

 


	\item (\textit{USAMO 2015}) Let $S = \{1, 2, ..., n\}$, where $n \ge 1$. Each of the $2^n$ subsets of $S$ is to be colored red or blue. (The subset itself is assigned a color and not its individual elements.) For any set $T \subseteq S$, we then write $f(T)$ for the number of subsets of T that are blue.


 Determine the number of colorings that satisfy the following condition: for any subsets $T_1$ and $T_2$ of $S$,

 \[f(T_1)f(T_2) = f(T_1 \cup T_2)f(T_1 \cap T_2).\]


 



\textbf{Solution}

Define function: $C(T)=1$ if the set T is colored blue, and $C(T)=0$ if $T$ is colored red.Define the $\text{Core} =\text{intersection of all } T \text{ where } C(T)=1$. 

 The empty set is denoted as $\varnothing$, $\cap$ denotes intersection, and $\cup$ denotes union. Let $S_n=\{n\}$ are one-element subsets.

 Let $m_{c_k}=\dbinom{m}{k} = \frac{m!}{k!(m-k)!}$ denote m choose k.

 
(Case I) $f(\varnothing)=1$. Then for distinct m and k, $f(S_m \cup S_k)=f(S_m)f(S_k)$, meaning only if $S_m$ and $S_k$ are both blue is their union blue. Namely $C(S_m \cup S_k)=C(S_m)C(S_k).$

 Similarly, for distinct $m,n,k$, $f(S_m \cup S_k \cup Sn)=f(S_m \cup S_k)f(S_n)$, $C(S_m \cup S_k \cup S_n)=C(S_m)C(S_k)C(S_n)$. This procedure of determination continues to $S$. Therefore, if $T=\{a_1,a_2, \cdots a_k\}$, then $C(T)=C(S_{a_1})C(S_{a_2}) \cdots C(S_{a_k})$. All colorings thus determined by the free colors chosen for subsets of one single elements satisfy the condition.  There are $2^n$ colorings in this case. 

 (Case II.)  $f(\varnothing)=0$.

 (Case II.1)  $\text{Core}=\varnothing$. Then either (II.1.1) there exist two nonintersecting subsets A and B, $C(A)=C(B)=1$, but f$(A)f(B)=0$, a contradiction, or (II.1.2) all subsets has $C(T)=0$, which is easily confirmed to satisfy the condition $f(T_1)f(T_2)=f(T_1 \cap T_2)f(T_1 \cup T_2)$.  There is one coloring in this case.

 (Case II.2) Core = a subset of 1 element. WLOG, $C(S_1)=1$. Then $f(S_1)=1$, and subsets containing element 1 may be colored blue. $f(S_1 \cup S_m)f(S_1\cup S_n)=f(S_1 \cup S_m \cup S_n)$, namely $C(S_1 \cup S_m \cup S_n)=C(S_m \cup S_1)C(S_n \cup S_1)$. Now S\_1 functions as the $\varnothing$ in case I, with $n-1$ elements to combine into a base of $n-1$ two-element sets, and all the other subsets are determined. There are $2^{n-1}$ colorings for each choice of core. However, there are nC1 = n such cores. Hence altogether there are $n2^{n-1}$ colorings in this case.

 (Case II.3) Core = a subset of 2 elements. WLOG, let $C(S_1 \cup S_2)=1$. Only subsets containing the core may be colored blue. With the same reasoning as in the preceding case, there are $(nC2)2^{n-2}$ colorings.

 $\dots$

 (Case II.n+1) Core = S. Then $C(S)=1$, with all other subsets $C(T)=0$, there is $1=\dbinom{n}{n}2^0$

 Combining all the cases, we have $1+\left[2^n+\dbinom{n}{1}2^{n-1}+\dbinom{n}{2}2^{n-2}+ \cdots + \dbinom{n}{n}2^0\right]=\boxed{1+3^n}$ colorings. sponsored by ALLEN

 


	\item (\textit{USAMO 2015}) Steve is piling $m\geq 1$ indistinguishable stones on the squares of an $n\times n$ grid. Each square can have an arbitrarily high pile of stones. After he finished piling his stones in some manner, he can then perform stone moves, defined as follows. Consider any four grid squares, which are corners of a rectangle, i.e. in positions $(i, k), (i, l), (j, k), (j, l)$ for some $1\leq i, j, k, l\leq n$, such that $i<j$ and $k<l$. A stone move consists of either removing one stone from each of $(i, k)$ and $(j, l)$ and moving them to $(i, l)$ and $(j, k)$ respectively, or removing one stone from each of $(i, l)$ and $(j, k)$ and moving them to $(i, k)$ and $(j, l)$ respectively.


 Two ways of piling the stones are equivalent if they can be obtained from one another by a sequence of stone moves.


 How many different non-equivalent ways can Steve pile the stones on the grid?


 



\textbf{Solution}

Let the number of stones in row $i$ be $r_i$ and let the number of stones in column $i$ be $c_i$. Since there are $m$ stones, we must have $\sum_{i=1}^n r_i=\sum_{i=1}^n c_i=m$

 Lemma 1: If any $2$ pilings are equivalent, then $r_i$ and $c_i$ are the same in both pilings $\forall i$.

 Proof: We suppose the contrary. Note that $r_i$ and $c_i$ remain invariant after each move, therefore, if any of the $r_i$ or $c_i$ are different, they will remain different.

 Lemma 2: Any $2$ pilings with the same $r_i$ and $c_i$ $\forall i$ are equivalent.

 Proof: Suppose piling 1 and piling 2 not the same piling. Call a stone in piling 1 wrong if the stone occupies a position such that there are more stones in that position in piling 1 than piling 2. Similarly define a wrong stone in piling 2. Let a wrong stone be at $(a, b)$ in piling 1. Since $c_b$ is the same for both pilings, we must have a wrong stone in piling 2 at column b, say at $(c, b)$, such that $c\not = a$. Similarly, we must have a wrong stone in piling 1 at row c, say at $(c, d)$ where $d \not = b$. Clearly, making the move $(a,b);(c,d) \implies (c,b);(a,d)$ in piling 1 decreases the number of wrong stones in piling 1. Therefore, the number of wrong stones in piling 1 must eventually be $0$ after a sequence of moves, so piling 1 and piling 2 are equivalent.

 Lemma 3: Given the sequences $g_i$ and $h_i$ such that $\sum_{i=1}^n g_i=\sum_{i=1}^n h_i=m$ and $g_i, h_i\geq 0 \forall i$, there is always a piling that satisfies $r_i=g_i$ and $c_i=h_i$ $\forall i$.

 Proof: We take the lowest $i$, $j$, such that $g_i, h_j >0$ and place a stone at $(i, j)$, then we subtract $g_i$ and $h_j$ by $1$ each, until $g_i$ and $h_i$ become $0$ $\forall i$, which will happen when $m$ stones are placed, because $\sum_{i=1}^n g_i$ and $\sum_{i=1}^n h_i$ are both initially $m$ and decrease by $1$ after each stone is placed. Note that in this process $r_i+g_i$ and $c_i+h_i$ remains invariant, thus, the final piling satisfies the conditions above.

 By the above lemmas, the number of ways to pile is simply the number of ways to choose the sequences $r_i$ and $c_i$ such that $\sum_{i=1}^n r_i=\sum_{i=1}^n c_i=m$ and $r_i, c_i \geq 0 \forall i$. By stars and bars, the number of ways is $\binom{n+m-1}{m}^{2}$.

 Solution by Shaddoll

 


	\item (\textit{USAMO 2015}) Let $a, b, c, d, e$ be distinct positive integers such that $a^4 + b^4 = c^4 + d^4 = e^5$. Show that $ac + bd$ is a composite number.


 



\textbf{Solution 1}

Note:  This solution is definitely not what the folks at MAA intended, but it works!

 Look at the statement $a^4+b^4=e^5$.  This can be viewed as a special case of Beal's Conjecture (\href{https://artofproblemsolving.comhttp://en.wikipedia.org/wiki/Beal\%27s\_conjecture}{http://en.wikipedia.org/wiki/Beal\%27s\_conjecture}), stating that the equation $A^x+B^y=C^z$ has no solutions over positive integers for $gcd(a, b, c) = 1$ and $x, y, z > 2$. This special case was proven in 2009 by Michael Bennet, Jordan Ellenberg, and Nathan Ng, as $(x, y, z) = (2, 4, n)$.  This case $a^4+b^4=e^5$ is obviously contained under that special case, so $a$ and $b$ must have a common factor greater than $1$.  

 Call the greatest common factor of $a$ and $b$ $f$.  Then $a = f \cdot a_1$ for some $a_1$ and likewise $b = f \cdot b_1$ for some $b_1$.  Then consider the quantity $ac+bd$.  

 $ac+bd = f \cdot a_1 \cdot c + f \cdot b_1 \cdot d = f\cdot(a_1 \cdot c + b_1 \cdot d)$.

 Because $c$ and $d$ are both positive, $(a_1 \cdot c + b_1 \cdot d) > 1$, and by definition $f > 1$, so $ac+bd$ is composite.  

 



\textbf{Solution 2}

A more conventional approach, using proof by contradiction:

 Without loss of generality, assume $a>d$ 

 Since $a^4 +b^4=c^4+d^4$, It is obvious that $b<c$

 We construct the equation $(a^4 +b^4)c^2d^2-(c^4+d^4)a^2b^2=(a^2c^2-b^2d^2)(a^2d^2-b^2c^2)$ (the right side is the factorization of the left) Which factors into: $e^5(cd-ab)(cd+ab)=(ac-bd)(ac+bd)(ad-bc)(ad+bc)$ (by using $a^4 +b^4=c^4+d^4=e^5$ and factoring)

 If $ac-bd$ or $ad-bc$ equal zero, then $a:b$ equals $c:d$ or $d:c$ respectively, which is impossible because $a^4 +b^4=c^4+d^4$ and because $a,b,c,d,e$ are distinct. Therefore the right side of the equation above is non-zero and the left side must be divisible by $ac+bd$ 

 If $ac+bd$ were prime,

 We have $ac+bd-(cd+ab)=(a-d)(c-b)>0$

 Which means that  $ac+bd>cd+ab>cd-ab$ and neither $cd+ab$ nor $cd-ab$ can be multiples of $ac+bd$ meaning $e$ must be a multiple of $ac+bd$ and $e=k(ac+bd)$ for some integer $k$

 But clearly this is impossible since $(k(ac+bd))^5>a^4+b^4$.

 Therefore, by contradiction, $ac+bd$ is composite.

 


	\item (\textit{USAMO 2015}) Consider $0<\lambda<1$, and let $A$ be a multiset of positive integers. Let $A_n=\{a\in A: a\leq n\}$. Assume that for every $n\in\mathbb{N}$, the set $A_n$ contains at most $n\lambda$ numbers. Show that there are infinitely many $n\in\mathbb{N}$ for which the sum of the elements in $A_n$ is at most $\frac{n(n+1)}{2}\lambda$. (A multiset is a set-like collection of elements in which order is ignored, but repetition of elements is allowed and multiplicity of elements is significant. For example, multisets $\{1, 2, 3\}$ and $\{2, 1, 3\}$ are equivalent, but $\{1, 1, 2, 3\}$ and $\{1, 2, 3\}$ differ.)


 



\textbf{Solution}

Proposed by mengmeng142857.

 We prove this by contradiction. Suppose that the number of times that an integer $a$ appears in $A$ is $k_a$. Let the sum of the elements in $A_n$ be $S_n=\sum_{i=1}^{n}i\cdot k_i$. If there are only finitely many $n\in\mathbb{N}$ such that $S_n\leq\frac{n(n+1)}{2}\lambda$, then by the Well Ordering Principle, there must be a largest $m\in\mathbb{N}$ such that $S_m\leq\frac{m(m+1)}{2}\lambda$, and for all $n>m$, $S_n>\frac{n(n+1)}{2}\lambda$.

 Now, observe that for some $n>m$,
 \[\frac{1}{n}\cdot S_{n}+\frac{1}{n(n-1)}\cdot S_{n-1}+\frac{1}{(n-1)(n-2)}\cdot S_{n-2}+...+\frac{1}{2\cdot1}\cdot S_{1} =\sum_{i=1}^{n}k_i \leq n\lambda\]Also, $\frac{S_{n}}{n}>\frac{n+1}{2}\lambda$, $\frac{S_{n-1}}{n(n-1)}>\frac{1}{2}\lambda$, $\frac{S_{n-2}}{(n-1)(n-2)}>\frac{1}{2}\lambda$, $\hdots$$\frac{S_{m+1}}{(m+2)(m+1)}>\frac{1}{2}\lambda$, $\frac{S_{m}}{(m+1)m}\leq \frac{1}{2}\lambda$.

 Now, for any $n\in\mathbb{N}$, we let $\frac{S_n}{(n+1)n}-\frac{1}{2}\lambda=d_n$, then
 \[\frac{1}{n}\cdot S_n+\sum_{i=1}^{n-1}\frac{S_i}{i(i+1)} =n\lambda+(n+1)\cdot d_n+\sum_{i=m+1}^{n-1}d_i+\sum_{i=1}^{m}d_i \leq n\lambda\]Let $\sum_{i=1}^{m}d_i=-c$ (where $c$ is positive, otherwise we already have a contradiction), then 
 \[n\lambda-c< n\lambda+(n+1)\cdot d_n+\sum_{i=m+1}^{n-1}d_i-c =\sum_{i=1}^{n}k_i\leq n\lambda\]Dividing both sides by $n$, $\lambda-c/n<\frac{\sum_{i=1}^{n}k_i}{n}\leq \lambda$. Taking the limit yields$\lim_{n\to\infty}{\frac{\sum_{i=1}^{n}k_i}{n}}=\lambda$.Observe that this would imply that
 \[\lim_{n\to\infty}{k_n} =\lim_{n\to\infty}{\left(\sum_{i=1}^{n}k_i-\sum_{i=1}^{n-1}k_i\right)} =\lim_{n\to\infty}{\left(n\cdot\frac{\sum_{i=1}^{n}k_i}{n}-(n-1)\cdot\frac{\sum_{i=1}^{n-1}k_i}{n-1}\right)} =\lambda\]However, as $k_n$ is an integer, it is impossible for it to converge to $\lambda$. This yields the desired contradiction.

 


\end{enumerate}